\subsection{Cinemática e relação de velocidades}
\label{sec:cinematica_velocidades}

Nesta subseção será estabelecida a relação cinemática entre a base $\{B\}$ e o efetuador $\{E\}$ do MR. Será evidenciado através do acoplamento por meio de um Jacobiano que inclui as velocidades generalizadas da base e das juntas. O referencial utilizado é o da base $\{B\}$.

A pose (orientação e posição) do efetuador é dada por \eqref{eq:transformacao_total}, enquanto que a posição do efetuador em relação à base $\{B\}$ é o vetor coluna de translação de \eqref{eq:transformacao_total}:

\begin{equation}
\label{eq:rBE}
{}^{B}\vec r_{BE} \;=\; \text{pos}\big({}^{B}T_{E}(\vec\theta)\big).
\end{equation}

Para definir as velocidades generalizadas da base $\{B\}$, faz-se:

\begin{equation}
\label{eq:nu_base}
\dot{\vec q}_B \;=\;
\begin{bmatrix}
{}^{B}\vec v_B \\[2pt]
{}^{B}\vec \omega_B
\end{bmatrix},
\qquad
\dot{\vec\theta} \;=\; [\,\dot\theta_1,\dot\theta_2,\dot\theta_3,\dot\theta_4,\dot d_5\,]^T.
\end{equation}

A velocidade espacial (twist) do efetuador $\{E\}$, também em relação à base $\{B\}$ é dado por:

\begin{equation}
\label{eq:twist_coupled}
\begin{bmatrix}
{}^{B}\vec v_E \\[2pt]
{}^{B}\vec \omega_E
\end{bmatrix}
\;=\;
J_B(\vec q)\,
\begin{bmatrix}
{}^{B}\vec v_B \\[2pt]
{}^{B}\vec \omega_B
\end{bmatrix}
\;+\;
J_m(\vec q)\,\dot{\vec\theta},
\end{equation}

onde $J_B$ é o Jacobiano que propaga a velocidade espacial da base até o efetuador; enquanto que $J_m$ é o Jacobiano geométrico do manipulador.

O termo $J_B$ decorre da cinemática de um corpo rígido:

\begin{equation}
\label{eq:JB}
J_B(\vec q)
\;=\;
\begin{bmatrix}
I_{3} & -\,[\,{}^{B}\vec r_{BE}\,]_{\times} \\
0_{3\times 3} & I_{3}
\end{bmatrix},
\end{equation}

de modo que:

\begin{equation}
\label{eq:vE_from_vB}
{}^{B}\vec v_E \;=\; {}^{B}\vec v_B \;+\; {}^{B}\vec \omega_B \times {}^{B}\vec r_{BE},
\qquad
{}^{B}\vec \omega_E \;=\; {}^{B}\vec \omega_B.
\end{equation}

O Jacobiano do manipulador é construído coluna a coluna, a partir dos eixos articulares em relação à base $\{B\}$.
Considerando ${}^{B}\hat z_i$ o eixo da junta $i$ e ${}^{B}\vec r_i$ a origem do referencial da junta $i$, ambos obtidos de ${}^{B}T_i$:

\begin{equation}
\label{eq:Jm_colunas}
J_m(\vec q) \;=\; \big[\, \mathcal{J}_1 \;\mathcal{J}_2 \;\mathcal{J}_3 \;\mathcal{J}_4 \;\mathcal{J}_5 \,\big],
\end{equation}

com

\begin{equation}
\label{eq:Jm_rev}
\mathcal{J}_i \;=\;
\begin{bmatrix}
{}^{B}\hat z_i \times \big({}^{B}\vec r_{BE} - {}^{B}\vec r_i\big) \\
{}^{B}\hat z_i
\end{bmatrix}
\quad \text{(junta revoluta)},
\qquad
\mathcal{J}_5 \;=\;
\begin{bmatrix}
{}^{B}\hat z_5 \\
0_{3\times 1}
\end{bmatrix}
\quad \text{(junta prismática)}.
\end{equation}

Para obter apenas a velocidade linear do efetuador, utiliza-se a parte translacional de \eqref{eq:twist_coupled}:

\begin{equation}
\label{eq:linear_relation}
{}^{B}\dot{\vec r}_{E}
\;=\;
{}^{B}\vec v_B
\;+\;
{}^{B}\vec \omega_B \times {}^{B}\vec r_{BE}
\;+\;
J_{m,\text{lin}}(\vec q)\,\dot{\vec\theta},
\end{equation}
onde $J_{m,\text{lin}}$ é a submatriz $3\times 5$ composta pelas três primeiras linhas de $J_m$.

Para converter as velocidades da base entre referenciais, usa-se a rotação ${}^{I}R_B$:

\begin{equation}
\label{eq:ref_map}
\begin{bmatrix}
{}^{I}\vec v_B \\[2pt]
{}^{I}\vec \omega_B
\end{bmatrix}
\;=\;
\begin{bmatrix}
{}^{I}R_B & 0 \\
0 & {}^{I}R_B
\end{bmatrix}
\begin{bmatrix}
{}^{B}\vec v_B \\[2pt]
{}^{B}\vec \omega_B
\end{bmatrix}.
\end{equation}

Para encontrar o centro de massa (CoM) total e as posições relativas, considera-se a massa $m_k$ e e a posição do centro de massa do corpo ${}^{B}\vec r_{G_k}$, sendo que $k \in \{B,1,\dots,5\}$, expressa em relação à base $\{B\}$:

\begin{equation}
\label{eq:cm_total}
M \;=\; \sum_{k} m_k,
\qquad
{}^{B}\vec r_{G} \;=\; \frac{1}{M}\sum_{k} m_k\,{}^{B}\vec r_{G_k}.
\end{equation}

As posições relativas dos elos seguem das transformações homogêneas:

\begin{equation}
\label{eq:ri_pos}
{}^{B}\vec r_i \;=\; \text{pos}\big({}^{B}T_{i}\big),
\qquad
{}^{B}\vec r_{G_i} \;=\; {}^{B}R_i\,{}^{i}\vec r_{G_i} \;+\; {}^{B}\vec r_i,
\end{equation}

onde ${}^{i}\vec r_{G_i}$ é a posição do centro de massa do elo $i$ no seu próprio referencial e ${}^{B}R_i$ é a rotação de $\{i\}$ para $\{B\}$ extraída de ${}^{B}T_i$.